\documentclass[12pt,nonacm,acmlarge,titlepage,screen]{acmart}

% make ACM settings and load extra packages
% settings specific to the ACM class

%% DO NOT MODIFY THIS FILE UNLESS IT IS FOR SUBMITTING TO AN ACM CONFERENCE/JOURNAL

%% Rights management information.  This information is sent to you
%% when you complete the rights form.  These commands have SAMPLE
%% values in them; it is your responsibility as an author to replace
%% the commands and values with those provided to you when you
%% complete the rights form.
% \setcopyright{acmlicensed}
\setcopyright{none}
% \copyrightyear{2018}
% \acmYear{2018}
\acmDOI{XXXXXXX.XXXXXXX}

%% These commands are for a JOURNAL article.
% \acmJournal{POMACS}
% \acmVolume{37}
% \acmNumber{4}
% \acmArticle{111}
% \acmMonth{8}

%% Submission ID.
%% Use this when submitting an article to a sponsored event. You'll
%% receive a unique submission ID from the organizers
%% of the event, and this ID should be used as the parameter to this command.
% \acmSubmissionID{123-A56-BU3}


%% citation style
% \citestyle{acmauthoryear}
\citestyle{acmnumeric}


%% ACM dates
% \received{\docDate}
% \received[revised]{12 March 2009}
% \received[accepted]{5 June 2009}
%% This file contains extra packages to use on top of the ACM template

%% STYLING ToC
\usepackage{titletoc}

%% SUB-FIGURES
\usepackage{subcaption}

%% ALGORITHMS
\usepackage{algorithm}
\usepackage{algorithmicx}
\usepackage{algpseudocode}

%% MORE EQUATIONS TOOLS
\usepackage{amsmath,mathtools}

%% CODE LISTINGS
\usepackage{listings}

% EASIER CROSS-REFERENCES
\usepackage[capitalise,nameinlink,noabbrev]{cleveref}

% TO GET ORDINAL NUMBERS
\usepackage[super]{nth}

% PAGE-BREAKABLE TABLES
\usepackage{longtable}

% IN-TABLE TEXT ALIGNMENTS
\usepackage{array}

\usepackage{enumitem} % requirements list
% \usepackage{tabularx} % for requirements table
% \usepackage{booktabs}

% document details
% replace those with your details
\def\docTitle{Hardware Acceleration for Skin Lesion Classification Using an Embedded FPGA}
\def\docAuthor{Lewis Turner}
\def\docAuthorID{H00430662}

% \def\docSupervisor{Dr. Rob \textsc{Stewart}, Ms. Tess \textsc{Watt}}
\newcommand{\docSupervisorOne}{Dr. Rob Stewart}
\newcommand{\docSupervisorTwo}{Ms. Tess Watt}
\newcommand{\docSupervisor}{% stacked names for title page
	\\
	\docSupervisorOne
	\\
	\docSupervisorTwo
}

\def\docDegree{BSc (Hons) Computer Science}
%\def\docDegree{MSc Cookery}
\def\docType{Research Report}
% \def\docType{Honours Dissertation}
% \def\docType{Masters Dissertation}
\def\docCourse{F20PA/B/C - Honours Project}
% \def\docCourse{F21RP - Research Methods and Project Planning}
% \def\docCourse{F21MP - Masters Project and Dissertation}

% load extra formatting
% formatting and package settings


% document details
\def\docDate{\nth{\the\day}~\ifcase\the\month\or
  January\or February\or March\or April\or May\or June\or
  July\or August\or September\or October\or November\or
  December\fi~\the\year}
\def\docInstitution{Heriot-Watt University}
\def\docSchool{School of Mathematical and Computer Sciences}
\author{\docAuthor}
\title{\docTitle}
\date{\docDate}

% change geometry of ACM template
\geometry{twoside=true, head=13pt, foot=2pc,
     paperwidth=21cm, paperheight=29.7cm,
     includeheadfoot,
     top=78pt, bottom=114pt, inner=81pt, outer=81pt,
     marginparwidth=4pc,heightrounded
 }%

% to make sections start on new pages
% \AddToHook{cmd/section/before}{\newpage}

% colors
% main text
\def\textColor{\color[HTML]{212121}} 
\def\secTextColor{\color[HTML]{424242}}
% footer and header
\def\headerColor{\color[HTML]{565656}} 
% title, chapters, sections
\def\titleColor{\color[HTML]{212121}}

% set line spacing
% \onehalfspacing
% \singlespacing

% set table cells padding
\def\arraystretch{1.5}

% set headings format
% unset to follow ACM format
% \titleformat{\chapter}{\singlespacing\normalfont\huge\bfseries\titleColor}{\chaptertitlename\ \thechapter}{1em}{}
% \titleformat{\chapter}{\singlespacing\normalfont\huge\bfseries\titleColor}{\thechapter}{1em}{}
% \titleformat{\section}{\singlespacing\normalfont\large\bfseries\titleColor}{\thesection}{1em}{}
% \titleformat{\subsection}{\singlespacing\normalfont\large\bfseries}{\thesubsection}{1em}{}
% \titleformat{\subsubsection}{\singlespacing\normalfont\normalsize\bfseries}{\thesubsubsection}{1em}{}
% \titlespacing*{\chapter}{0pt}{0pt}{10pt}

% set paragraph indent
% unset to follow ACM format
% \setlength{\parindent}{18pt} 

% set the text colour
% \textColor 

% set hyperlinks + pdf metadata
\hypersetup{
    pdftitle={\docTitle},
    pdfauthor={\docAuthor},
    pdfkeywords={computer science dissertation},
    colorlinks=true,%
    citecolor=[HTML]{152845},%
    linkcolor=[HTML]{152845},%
    urlcolor=[HTML]{1976D2}
}

% styling the code listings
\definecolor{codegreen}{rgb}{0,0.6,0}
\definecolor{codegray}{rgb}{0.5,0.5,0.5}
\definecolor{codepurple}{rgb}{0.58,0,0.82}
\definecolor{backcolour}{rgb}{0.95,0.95,0.92}

\lstdefinestyle{mystyle}{
    backgroundcolor=\color{backcolour},   
    commentstyle=\color{codegreen},
    keywordstyle=\color{magenta},
    numberstyle=\tiny\color{codegray},
    stringstyle=\color{codepurple},
    basicstyle=\ttfamily\singlespacing,
    breakatwhitespace=false,         
    breaklines=true,                 
    captionpos=b,                    
    keepspaces=true,                 
    numbers=left,                    
    numbersep=5pt,                  
    showspaces=false,                
    showstringspaces=false,
    showtabs=false,                  
    tabsize=2
}
\lstset{style=mystyle}

% reset names for code listing
\renewcommand{\lstlistingname}{Code Listing}
\renewcommand{\lstlistlistingname}{List of \lstlistingname s}

% set cref for listing
\crefname{lstlisting}{\lstlistingname}{\lstlistingname s}
\Crefname{lstlisting}{\lstlistingname}{\lstlistingname s}

% set cross-references labels with cleveref
\crefname{appendixchapter}{Appendix}{Appendices}
\Crefname{appendixchapter}{Appendix}{Appendices}


% set custom headers/footers
% unset to follow ACM template
% \pagestyle{fancy}
% \fancyhf{}
% \renewcommand{\headrulewidth}{0pt}
% \renewcommand{\footrulewidth}{0pt}
% \fancyhead[R]{\headerColor \docType}
% \fancyhead[L]{\headerColor \docAuthor}
% \fancyfoot[C]{\headerColor \thepage}
% % Redefine the plain page style for chapter pages
% \fancypagestyle{plain}{%
%     \fancyhf{}%
%     \fancyfoot[C]{\headerColor \thepage}
% }

% custom section heading for front-matter content with addition to the ToC
\newcommand{\frontmattersection}[2]{
    \thispagestyle{plain} % remove the header on page
    % \vspace*{1cm} % to align vertically
    \phantomsection
    \begin{center} 
        \textsc{#1} 
        \addcontentsline{toc}{section}{#1}
    \end{center}
    #2 % Main Text
    \clearpage % leave blank space underneath
}

\renewcommand{\abstract}[1]{
    \frontmattersection{Abstract}{#1}
}

\newcommand{\acknowledgements}[1]{
    \frontmattersection{Acknowledgements}{#1}
}

% formatting the table of content
\dottedcontents{section}[1em]{\bfseries}{1.2em}{0pc}
\dottedcontents{subsection}[3em]{}{2.2em}{1pc}

\renewcommand{\contentsname}{Table of Contents}

% objectives list formatting
\newlist{objs}{enumerate}{1}
\setlist[objs]{
	label=\textbf{O\arabic*}, % appearance in list
	ref=[O\arabic*] % appearance when referencing
}

% functional requirements list formatting
\newlist{reqs}{enumerate}{1}
\setlist[reqs]{
	label=\textbf{FR\arabic*}, % appearance in list
	ref=[FR\arabic*], % appearance when referencing
	% leftmargin=*,
	% align=left
}
% \crefname{reqsi}{Requirement}{reqs} % ref link appearance of "Requirement [R?]"

% non-function requirments list formatting
\newlist{reqsn}{enumerate}{1}
\setlist[reqsn]{
	label=\textbf{NR\arabic*}, % appearance in list
	ref=[NR\arabic*] % appearance when referencing
}

% custom \cref which surrounds link with square brackets (% below is necessary to stop extra spaces appearing)
% name being bcref because VSCode ref syntax looks for -ref command
% s - if added * to end to not create hyperlink
% m - mandatory label inside {}
\NewDocumentCommand{\bcref}{s m}{%
	[\IfBooleanTF{#1}{\cref*{#2}}{\cref{#2}}]%
}

\begin{document}

%%%%%%%%%%%%%%%%%%%%%%%%%%%%%%%%%%%%%%%%%%%%%%%%%%
% FRONT MATTER
%%%%%%%%%%%%%%%%%%%%%%%%%%%%%%%%%%%%%%%%%%%%%%%%%%
\begin{titlepage}

\begin{center}
	
	% Title (and subtitle)
	~\\[0.2cm]
	{\huge \titleColor \textbf{\docTitle}}\\
	~\\[0.7cm]
	
	% Author(s)
	{\Large \textit{by} \docAuthor}\\[0.7cm]
	{\large \docType}\\[1cm]
	{\large \textit{Submitted in partial fulfilment for the degree of}\\[0.2cm] \docDegree}\\[0.7cm]
	{\large \textit{Supervised by} {\Large \docSupervisor}}
	~\\[1cm]
	
	% Organisation	
	% ~\\[0cm]
        % logo / coat of arms
        \begin{figure}[h]
            \centering
		\includegraphics[height=3cm,keepaspectratio]{images/hwcoa_bw.jpg}
        \Description{Heriot-Watt University's coat-of-arms}
	\end{figure}
	\begin{large}
		\textsc{\docInstitution}\\
		\docSchool\\
% 		Department\\
	\end{large}
	~\\[1cm]
	
	% Date
	{\large \docDate}\\
	~\\[1.5cm]
	
	% Copyright
	\begin{minipage}{0.7\textwidth}
	% \singlespacing
	\begin{footnotesize}
	The copyright in this dissertation is owned by the author. Any quotation from the dissertation or use of any of the information contained in it must acknowledge it as the source of the quotation or information.
	\end{footnotesize}
	\end{minipage}
	
\end{center}

\end{titlepage}

\newpage
\pagenumbering{roman}

% authorship declaration
% this declaration is automatically filled with the details provided in main.tex

\thispagestyle{plain} % remove the header on page

\phantomsection
\begin{center} 
	\textsc{Declaration of Authorship} 
	\addcontentsline{toc}{section}{Declaration of Authorship}
\end{center}


\begin{center} 
	\begin{tabular}{|p{0.3\textwidth}|p{0.675\textwidth}|}
		\hline
		\textbf{Course code and name} & \docCourse \\
		\hline
		\textbf{Type of assessment} & Individual \\
		\hline
		\textbf{Document title} & \docTitle \\
		\hline
		\textbf{Student name} & \docAuthor \\
		\hline
		\textbf{Student ID} & \docAuthorID \\
		\hline
	\end{tabular}
\end{center}

\bigskip

\begin{longtable}{|>{\raggedright\arraybackslash}p{\textwidth}|}
	\hline
	\endfirsthead
	\hline
	\multicolumn{1}{|r|}{\footnotesize\it (form continued from previous page)} \\
	\endhead
	\multicolumn{1}{|r|}{\footnotesize\it (form continues on next page)} \\
	\hline
	\endfoot
	\hline
	\endlastfoot
	
		% EDIT this statement to indicate what use you made of Generative AI for your project work and in preparation of this document. If your project is about developing or experimenting with Generative AI, you still need to indicate your Generative AI use beyond the developments and experiments. If you have not used Generative AI, indicate it explicitly.
		\textbf{Statement on Generative AI Use:}
		\\
		Several updates to Google Gemini models came out during the writing of this document, the wording used below matches features available at the time from https://gemini.google.com/app with a personal Gemini Pro family account.
		\\
		\begin{itemize}
			\item Google Gemini 4.0 Argon (Thinking level: Low - High) to assist debugging while setting up FPGA for use over Wifi
			\item Google Gemini 3.8 Flash (Extended Thinking) and Gemini 4 Hydrogen (Thinking level: Low - High) to help search for relevant research papers. We would ask it to create a table consisting of papers: title, year, summary, DOI. This allowed best leverage of LLM's speed for searching the internet like a more powerful keyword search.
			\item Gemini 4 Hydrogen (Thinking level: High) to assist identifying relevant and compatible tools for the methodology pipeline.
		\end{itemize}
		\\
	\hline
	\multicolumn{1}{l}{}
	\\
	\hline
		\textbf{By signing this form:}
		\\
		
		\textbullet~ I declare that the work I have submitted for this assessment is my own original work. I have not presented the ideas, words, data, code, designs, analyses, or other intellectual contributions of another person or system as my own. 
		\\
		
		\textbullet~ My submission is expressed in my own words except where the work of others has been appropriately acknowledged, cited, and referenced. All sources used in the preparation of this assessment, including books, journal articles, websites, software tools, and Generative AI tools, have been properly acknowledged.
		\\
		
		\textbullet~ Where Generative AI tools have been used, I have complied with the course’s stated policies regarding Generative AI use. I have clearly and accurately disclosed in the above statement box all uses of Generative AI in the assessment, including the specific tools used, the purpose for which they were used, and the extent to which any generated content influenced the final submission.
		\\
		
		\textbullet~ I understand that I remain fully responsible for the content, accuracy, integrity, and originality of my submission, regardless of any assistance received from Generative AI or other tools.
		\\

		\textbullet~ I confirm that I have read, understood and followed the University's Regulations on academic misconduct as published on the \href{https://www.hw.ac.uk/uk/services/academic-registry/academic-integrity/academic-misconduct.htm}{University's website}, and that I am aware of the penalties that I will face should I not adhere to the University Regulations.
		\\
		
		\textbullet~ I confirm that I have read, understood and avoided the different types of plagiarism explained in the University guidance on \href{https://heriotwatt.sharepoint.com/sites/skillshub/SitePages/Academic-Integrity-and-Plagiarism.aspx}{Academic Integrity and Plagiarism}.
		\\
		
		\\

		\noindent \textbf{Student signature}: \docAuthor
		\\
		
		\\
		
		\noindent \textbf{Date}: \docDate
		\\
	\hline
\end{longtable}

% \noindent \textbf{Your work will not be marked if a signed copy of this form is not included with your submission.}



\clearpage

% Acknowledgements
\acknowledgements{
Acknowledgements text.
}

% Abstract
\abstract{
Abstract text.
}

% table of contents

\newpage

\thispagestyle{plain} % remove the header on page

\setcounter{tocdepth}{2} % set depth of ToC, 0 = cha, 1 = sec, etc.

\tableofcontents % the actual table of contents command

\clearpage

	
%%%%%%%%%%%%%%%%%%%%%%%%%%%%%%%%%%%%%%%%%%%%%%%%%%
% MAIN BODY
%%%%%%%%%%%%%%%%%%%%%%%%%%%%%%%%%%%%%%%%%%%%%%%%%%
\newpage
\pagenumbering{arabic}

% Introduction
%%%%%%%%%%%%%%%%%%%%
\section{Introduction} \label{cha:intro}
Summary of project motivation \& context

	\subsection{Background and Motivation} \label{sec:intro_motiv}
	\textbullet~ Using papers (mainly) from literature review to provide background context\\
	\textbullet~ Motivations for doing this project\\
	\textbullet~ Identifying problems to be solved to achieve objectives

	\subsection{Aim and Objectives} \label{sec:intro_aim}
	State project aims\\
	Table of objectives with IDs to reference later

	\begin{objs}
		\item\label{obj:train} Train/download a skin lesion model
		\item\label{obj:compress} Compress skin lesion model successfully for FPGA
		\item\label{obj:run} Run compressed model on FPGA successfully
	\end{objs}

	Testing referring to O - \ref{obj:run}

% Plan of the document.

% Background Research
%%%%%%%%%%%%%%%%%%%%

% RADAR: Relevance, Authority, Date, Appearance, Reason for writing
%
%	- paragraph summary on each paper/topic
%	- paragraphs cross reference each other
%
% 1. starting with finding current skin lesion models
% 2. choosing mobile hardware as target: mobile, RPi, FPGA (biggest part)
% 3. find pipelines for using FPGA
% 4. compression techniques (assuming model too big at baseline)
\newpage
\section{Literature Review} \label{sec:lit_rev}
This section covers my research for the areas of this project. First finding generally what a skin lesion classification model is, then looking into recent models \bcref{lit:models}. Then how I chose an FPGA as my target hardware \bcref{lit:hardware}. What methods exist for running a ML model on an FPGA \bcref{lit:fpga_pipelines}. Recent image classifier compression techniques so the model running on the FPGA has an accuracy closer to the baseline model \bcref{lit:compression}. The methods I used for finding papers are in \bcref{lit:search_methods}.

	% \subsection{Acknowledgements}
	% Sources for literature/statements
	% \\
	% Continuing idea by supervisor Tess Watt in her paper but deployed to FPGA instead of RPi \cite{tess_pi}
	% \\
	% Survey of recent methods for skin lesion classification between 2017 and 2023 \cite{survey_2023}

	% How I searched for papers, what tools did I use and how did I verify my sources
	\subsection{Search Methods} \label{lit:search_methods}
	I used a combination of Google Gemini and Google Scholar to do keyword searches for relevant papers. For all papers I started with reading the abstract then went to other sections if I required more information.

	% Checklist
	% - minimum 5 papers unrelated to FPGAs on skin lesions
	% - paragraph summary on each paper/topic
	% - paragraphs cross reference each other
	% 	
	% Papers on recent skin classification models
	\subsection{Skin Lesion Classification Models} \label{lit:models}
	Recent classification models use a variety of datasets that each have varying types of identified skin lesions (cite a few papers that use different datasets). This subsection covers commonly used datasets and looks at some recent models with high accuracy.

		\subsubsection{Datasets}
		There are a few popular datasets used to train these models.
		\\
		Ham10000 contains 10,015 dermoscopic images intended for academic machine learning use. It has 7 different categories of skin lesions which is a "representative collection of all important diagnostic categories in the realm of pigmented lesions" \cite{ham_src}.
		\medskip

		The ISIC Archive has several versions and challenge datasets available. The ISIC 2020 version contains 23,385 images with 8 classes of identified skin lesions and 1 class for unknown lesions.
		\medskip

		Other datasets exist such as PH2 \cite{ph2_paper} which are designed to identify a specific type of lesion but these datasets tend to be much smaller and can't effectively be used for generalised classification.

		\subsubsection{Recent Models Performance}
		They use however different methods and techniques with varying performance. Performance is typically measured by model accuracy for the dataset it was trained on.
		\\
		Results of models are typically compared in terms of classification accuracy on a test dataset unseen to the models \cite{survey_2023} (cite both years surveys). A model is typically considered better if the accuracy is as close to 100\% as possible without the model overfitting. An overfitting model would appear to have high accuracy during testing but then have lower accuracy on unseen data (cite overfitting definition).

		\subsubsection{Model Architectures}
		Recent model architectures can achieve accuracies of 9X\% (cite models) which is a XX\% uplift from earlier fundamental papers such as MobileNet from 2017 \cite{mobilenet_2017}.
		\\
		There is no single model architecture, however I found a few of best performing ones from the last few years to compare:
		\begin{itemize}
			\item architecture 1 with benefits
			\item architecture 2 with benefits
			\item ...
		\end{itemize}

	% Looking into hardware options (RPi, mobile, FPGA)
	\subsection{Mobile Hardware for Machine Learning Acceleration} \label{lit:hardware}
	Many options for portable hardware exist such as mobile, Nvidia AI accelerators, RPi, embedded FPGAs. This section covers a few popular options including the PYNQ-Z1 I will be using.
	
		\subsubsection{Raspberry Pi 3}
		RPi boards are designed as an embedded System on Chip (SoC) architecture however possess no specialised acceleration hardware beyond what is expected of a modern CPU. I am already familiar with the Raspberry Pi 3 \cite{pi3_product_page}. A cheap (around £30 at time of research) hardware package often used for running lightweight servers and connecting to a vast array of peripherals via a custom Linux operating system. Not known for high performance but has low power usage of only 1.5-4W according to \cite[Section 6]{tess_pi}.

		\subsubsection{Mobile}
		Modern mobile phones have various sensors, cameras and processors designed for general purpose use. For this project only the phone's processor and any AI accelerators would be used so many of the features would be unutilised. This amounts to a larger cost to have access to processors that exist on cheaper and simpler embedded SoC boards.

		\subsubsection{PYNQ-Z1}
		The PYNQ-Z1 is part of a series of boards by Xilinx (now owned by AMD) designed to make running neural networks on FPGAs more easily accessible \cite{pynq_z1_man}. It is an embedded SoC architecture that has an ARM CPU which can run a Linux bootloader similar to a RPi and an embedded Field Programmable Gate Array (FPGA) chip which can be used to create specialised accelerator hardware. Both chips share system resources including RAM \cite[Section 2]{pynq_z1_man}. An FPGA can achieve higher speeds than most accelerators while being programmable, only requiring a few Watts and having limited hardware resources. The PYNQ-Z1 provides a base framework to allow using the FPGA without having to design physical logic circuits. Many frameworks exist to generate efficient and optimised logic circuits which is explored in \bcref{lit:fpga_pipelines}.

	% Papers more generally on methodology for putting ML models on FPGA
	\subsection{Machine Learning Models on FPGA} \label{lit:fpga_pipelines}
	This section describes a few of the options for pipelines that could be used to run a machine learning model on an FPGA.

		% Mentioned as an option but only an old version supports PYNQ-Z1
		\subsubsection{Vitis AI}
		Vitis is a software platform used for digital signal processing and machine learning (cite vitis). The Vitis AI framework is part of the platform and is designed for AI acceleration of CNN models. A variety of tools and optimisations are included for designing programmable logic focusing on better area compute density and reducing power consumption compared to traditional implementations. It abstracts the programming to C/C++ languages level to allow for easier porting between specific hardware. The current version of Vitis AI is 5.0. The Zynq-7000 family (includes PYNQ-Z1) lost support as of version 1.4 because they were considered to not have enough resources for acceptable performance. Using this framework may still be possible but would be unlikely to have the same latest features as frameworks since 2021 (cite support).
		\\
		\cite[Page 14, AI Engine Development Environment]{vitis_paper}
		\\
		\cite[AMD Support Page, GitHub Vitis 1.4 Release]{vitis_1_4_support, vitis_1_4_git}

		% Explain here details about FINN framework: general process, benefits, drawbacks, supported model types
		\subsubsection{FINN}
		\cite{finn_paper}
		Preface with the FINN project and FINN compiler are separate but closely related projects. The FINN project uses the FINN compiler as part of a bigger pipeline.
		\\
		FINN is a pipeline for running quantised neural networks on FPGAs. It implements NN layers as streaming components meaning there are no processor clock cycles but instead the code is run directly in hardware. This means the NN runs in essentially constant time and with sub-microsecond latency. The main trade-off is that 
		- Runs NN in constant time determined by layer structure (execution completely input dependant)
		- Run speed variance is because of external factors

		- Designed for running image classifiers
		- Custom resources allocation for individual layers
		\cite{finn_project}
		

	% compression methods that best retain baseline model performance and best utilise FPGA
	\subsection{Compressing Machine Learning Models to FPGA} \label{lit:compression}
	Papers on useful compression techniques for certain model types
	% PAPERS ON COMPRESSION TECHNIQUES FOR ABOVE ARCHITECTURES

	\subsection{Further Research Recommendations}
	Identify gaps, need to repeat any research

	Gap I want to address: is it possible to take a skin lesion classification model unrelated to FPGAs and fit it entirely within PYNQ-Z1 FPGA hardware without significant accuracy loss.

\section{Methodology} \label{sec:method}
Detailed description of necessary methods, requirements, research questions that will be employed/
needed to complete the project

	\subsection{Plan for Execution}
	Explain how I plan to use experiments to fulfil all the requirements \& objectives

	\subsection{Experiment Planning}
	How I plan to conduct experiment(s) and what research question does it address/answer (detailed)

	% MOSCOW, enumerate so can be referenced later
	\subsection{Requirements}
	Quick explanation of functional requirements
	
	\begin{reqs}
		\item\label{req:model} Something about the skin lesion model
		\item\label{req:compress} Something about the compression
		\item\label{req:fpga} Something about the FPGA
	\end{reqs}
	
	Testing referring to FR - \ref{req:model}
	\\
	Quick explanation of non-function requirements

	\begin{reqsn}
		\item\label{reqn:speed} Something about running speed
		\item\label{reqn:compile_speed} Something about time taken to compile
	\end{reqsn}

	Testing referring to NR - \ref{reqn:compile_speed}

\section{Feasibility} \label{sec:feasiblity}

	\subsection{Preliminary Work}
	Summary of work carried out to prepare for this project to validate feasibility.

	\subsection{Evaluation Strategy}
	Identify for the project:
	\\
		\textbullet~ What needs evaluating
		\\
		\textbullet~ Why it should be evaluated
		\\
		\textbullet~ How it should be evaluated
		\\

\section{Conclusions}
Short summary of this doc

% Rest of the document...
%%%%%%%%%%%%%%%%%%%%


%%%%%%%%%%%%%%%%%%%%%%%%%%%%%%%%%%%%%%%%%%%%%%%%%%
% BACK MATTER
%%%%%%%%%%%%%%%%%%%%%%%%%%%%%%%%%%%%%%%%%%%%%%%%%%

% References
% list of references

\newpage
\renewcommand{\bibname}{References} % change name to References
% \phantomsection
% \addcontentsline{toc}{section}{References}
\bibliography{references}
% \bibliographystyle{plainnat}

% using Harvard style
% \bibliographystyle{agsm}

% using ACM style (defined in acm_settings.tex)
\bibliographystyle{ACM-Reference-Format}

% using APA style (requires apacite package)
% \bibliographystyle{apacite}

\section{Project Management}

	\subsection{Gantt Chart}
	Plan of activities to be carried out in semester 1, then future plan for semester 2

	\subsection{Analysis of Risks}
	Including mitigation plans

	\subsection{PLES}
	A well-researched consideration of any Professional, Legal, Ethical, and Social Issues
	pertinent to the project. (e.g. codes of conduct (BCS), codes of practice, standards,
	computer law, ethical decision making, intellectual property, social aspects, copyright,
	data protection, and so on).

% Summarising use of GenAI in project/deliverables
\section{Gen AI Summary} \label{sec:genai}
Searching for research papers I used Google Gemini 3.8 Flash on thinking mode High (October 2026 version) to source the most recent academic papers on a particular topic. I would ask it to create a table consisting of the papers: title, year, summary, DOI or similar URL. This allowed me to leverage the LLM's speed for searching the internet like a more powerful keyword search.

% Appendices
\newpage
\appendix

\crefalias{chapter}{appendixchapter} % for appendix crossreferences

% Rest of appendices...

\end{document}
